%! Linear Regression — Unit 1, Lesson 1.5 — Group Activity
\documentclass[10pt]{article}
\usepackage{algebra2-article}
\usepackage{algebra2-boxes}
\newcommand{\TallMath}[1]{$\displaystyle #1\rule[-1.4em]{0pt}{3.2em}$}
% Coordinate grid: {xmin}{xmax}{ymin}{ymax}
\newcommand{\lgrid}[4]{%
  \draw[step=1, linegray!40, very thin] (#1,#3) grid (#2,#4);
  \draw[->, semithick, charcoal] (#1,0)--(#2,0) node[below right, font=\scriptsize]{$x$};
  \draw[->, semithick, charcoal] (0,#3)--(0,#4) node[left, font=\scriptsize]{$y$};}
\newcommand{\dpt}[2]{\fill[navy] (#1,#2) circle (2.8pt);}

\begin{document}

\pageheader{Unit 1, Lesson 1.5}{Group Activity}

\vspace{0.08in}

\begin{headlinebox}{forestbg}
\small\textbf{Reading the Trend.} Real data scatters --- we summarize the \emph{trend}, not an exact rule.
Work with your partner. Describe an association by \textbf{direction} (up / down / none), \textbf{form}
(linear?), and \textbf{strength}. The \textbf{correlation coefficient} $r\in[-1,1]$ packs direction (sign)
and strength ($|r|$) into one number. The \textbf{line of best fit} $\hat y=ax+b$ lets you \emph{predict}
by substitution.
\end{headlinebox}

\vspace{0.06in}

\begin{tcolorbox}[colback=white, colframe=black!40, title=\textbf{Tier R --- Remediate},
    fonttitle=\bfseries, arc=1.5mm, left=3mm, right=3mm, top=2mm, bottom=2mm, breakable]
\begin{enumerate}[leftmargin=1.7em, itemsep=9pt]
  \item \textbf{Classify the direction} of each scatterplot: positive, negative, or none.
        \begin{center}
        \begin{tikzpicture}[scale=0.4]
          \lgrid{-0.5}{7}{-0.5}{7}
          \dpt{1}{1}\dpt{2}{3}\dpt{3}{3}\dpt{4}{5}\dpt{5}{6}
          \node[font=\scriptsize\bfseries] at (3,-1.6){(I)};
        \end{tikzpicture}
        \hspace{0.5cm}
        \begin{tikzpicture}[scale=0.4]
          \lgrid{-0.5}{7}{-0.5}{7}
          \dpt{1}{6}\dpt{2}{5}\dpt{3}{3}\dpt{4}{3}\dpt{5}{1}
          \node[font=\scriptsize\bfseries] at (3,-1.6){(II)};
        \end{tikzpicture}
        \end{center}
        (I) is \blank{2.4cm}; \quad (II) is \blank{2.4cm}.
  \item \textbf{Read a point.} In plot (I) above, the point farthest to the \emph{right} has coordinates
        $(\blank{0.8cm},\,\blank{0.8cm})$.
  \item \textbf{Predict with a given line.} A line of best fit is $\hat y=3x+5$.
        \\[3pt]
        $\hat y$ at $x=4$: \blank{1.6cm} \qquad $\hat y$ at $x=10$: \blank{1.6cm}
\end{enumerate}
\end{tcolorbox}

\begin{tcolorbox}[colback=white, colframe=black!40, title=\textbf{Tier A --- Approaching Proficiency},
    fonttitle=\bfseries, arc=1.5mm, left=3mm, right=3mm, top=2mm, bottom=2mm, breakable]
\begin{enumerate}[leftmargin=1.7em, itemsep=10pt]
  \item \textbf{Match $r$ to the plot.} Assign each value $r\in\{\,0.95,\ 0.10,\ -0.70\,\}$ to a plot.
        \begin{center}
        \begin{tikzpicture}[scale=0.4]
          \lgrid{-0.5}{6}{-0.5}{7}
          \dpt{1}{1}\dpt{2}{2}\dpt{3}{4}\dpt{4}{5}\dpt{5}{6}
          \node[font=\scriptsize\bfseries] at (2.5,-1.6){(P)};
        \end{tikzpicture}
        \hspace{0.4cm}
        \begin{tikzpicture}[scale=0.4]
          \lgrid{-0.5}{6}{-0.5}{7}
          \dpt{1}{4}\dpt{2}{2}\dpt{3}{5}\dpt{4}{3}\dpt{5}{4}
          \node[font=\scriptsize\bfseries] at (2.5,-1.6){(Q)};
        \end{tikzpicture}
        \hspace{0.4cm}
        \begin{tikzpicture}[scale=0.4]
          \lgrid{-0.5}{6}{-0.5}{7}
          \dpt{1}{5}\dpt{2}{3}\dpt{3}{4}\dpt{4}{2}\dpt{5}{2}
          \node[font=\scriptsize\bfseries] at (2.5,-1.6){(R)};
        \end{tikzpicture}
        \end{center}
        (P): $r=\blank{1.4cm}$; \quad (Q): $r=\blank{1.4cm}$; \quad (R): $r=\blank{1.4cm}$.
  \item \textbf{Interpret the model.} Relating \emph{years of experience} ($x$) to \emph{salary in \$1000s}
        ($y$), technology gives $\hat y=2.5x+38$ with $r=0.91$.
        \begin{enumerate}[label=(\alph*), leftmargin=1.8em, itemsep=4pt]
          \item The slope $2.5$ means: each additional year of experience is associated with about
                \blank{2.2cm} more in salary. \quad (Units? \blank{2.2cm})
          \item The $y$-intercept $38$ means the predicted starting salary (at $x=0$) is \blank{2.2cm}.
          \item Predict the salary at $x=10$ years.
                \begin{work}
                  \hat y &= 2.5(10) + 38 \\
                         &= 25 + 38      \\
                         &= 63
                \end{work}
                That is \$\blank{1.6cm}.
        \end{enumerate}
  \item \textbf{Interpolation or extrapolation?} The salaries above came from workers with $0$ to $15$ years
        of experience. Label each prediction:
        \\[3pt]
        $x=10$ years: \blank{2.6cm}; \qquad $x=30$ years: \blank{2.6cm}.
\end{enumerate}
\end{tcolorbox}

\begin{tcolorbox}[colback=white, colframe=black!40, title=\textbf{Tier E --- Extension},
    fonttitle=\bfseries, arc=1.5mm, left=3mm, right=3mm, top=2mm, bottom=2mm, breakable]
\begin{enumerate}[leftmargin=1.7em, itemsep=10pt]
  \item \textbf{Model a real trend.} A cafe records the daily high \emph{temperature} ($x$, in $^\circ$F)
        and \emph{cups of hot cocoa sold} ($y$). Over days from $20^\circ$ to $60^\circ$, technology gives
        \[ \hat y=-3x+260, \qquad r=-0.93. \]
        \begin{enumerate}[label=(\alph*), leftmargin=1.8em, itemsep=5pt]
          \item Describe the association (direction and strength), and say what the slope $-3$ means in
                context. \writelines{3}
          \item Predict cocoa sold at $x=40^\circ$: \blank{1.8cm}. Is this interpolation or extrapolation?
                \blank{2.4cm}
          \item Predict at $x=90^\circ$: \blank{1.8cm}. Explain why this prediction is \emph{unreasonable}.
                \writelines{2}
        \end{enumerate}
  \item \textbf{Correlation vs.\ causation.} A news story reports a \emph{strong positive} correlation
        between the number of firefighters sent to a fire and the dollar amount of damage, and concludes
        ``sending more firefighters causes more damage.''
        \begin{enumerate}[label=(\alph*), leftmargin=1.8em, itemsep=5pt]
          \item Is the conclusion valid? \blank{1.4cm} \; Name a \textbf{lurking variable} that explains the
                correlation. \writelines{2}
          \item In one sentence, state the difference between correlation and causation. \writelines{3}
        \end{enumerate}
\end{enumerate}
\end{tcolorbox}

\end{document}
